\section{A uniform sub Gaussian bound}\label{ssm:sec:binomial}
The eventual summation weight grows like the exponential of a positive square. Weak convergence alone cannot justify its expectation. We first establish a uniform bound strong enough to do so.

\begin{lemma}[Capped binomial exponential moments]\label{ssm:lem:subgaussian}
Let $X$ have the distribution of $\Bin(b,r)$ conditioned on $X\le k$, where $0<r<1$ and $0\le k\le b$ are arbitrary. Then
\begin{equation}\label{ssm:eq:variance}
 \Var(X)\le\frac b4,
\end{equation}
and, for every real $u$,
\begin{equation}\label{ssm:eq:mgf}
 \E e^{u(X-\E X)}\le e^{bu^2/8}.
\end{equation}
The same statements hold without conditioning.
\end{lemma}
\begin{proof}
Write the binomial as the sum of $b$ independent Bernoulli variables and condition their sum to be at most $k$. The conditional coordinates are exchangeable. For $b\ge2$, let $F(j)=\PP\{\Bin(b-2,r)\le j\}$, extended by zero below its support and by one above it. The unnormalized probabilities for two selected coordinates are
\[
 (1-r)^2F(k),\quad r(1-r)F(k-1),\quad
 r(1-r)F(k-1),\quad r^2F(k-2).
\]
Their covariance has the sign of $F(k-2)F(k)-F(k-1)^2$.

To check this sign directly, let $q_j$ be the binomial mass sequence, extended by zero. Its successive positive ratios decrease. Consequently, for $i\le j$,
$q_{j+1}q_{i-1}\le q_jq_i$. Sum over $i\le j$ to obtain
$q_{j+1}F(j-1)\le q_jF(j)$. Since
\[
 F(j)^2-F(j-1)F(j+1)=q_jF(j)-q_{j+1}F(j-1),
\]
$F$ is log-concave, including the support boundaries. Thus the conditional coordinates are pairwise negatively correlated. If $a$ is their common conditional success probability, then
\[
 \Var(X)\le b a(1-a)\le b/4.
\]
The cases $b\le1$ are immediate.

Exponential tilting by $e^{uX}$ preserves the capped-binomial family, replacing $r$ by
\[
 r_u=\frac{re^u}{1-r+re^u}.
\]
The second derivative of its log moment-generating function is the variance under this tilted law, hence is bounded by $b/4$ for every real $u$. Integrating twice from zero proves \eqref{ssm:eq:mgf}. Without a cap, the same variance and tilting argument applies.
\end{proof}

\begin{corollary}\label{ssm:cor:squareUI}
If $X_1,\ldots,X_M$ are independent, each ordinary or capped $\Bin(M-1,r)$ variables, set
\[
 Z_M=\frac{\sum_iX_i-\sum_i\E X_i}{M}.
\]
For $M\ge2$,
\begin{equation}\label{ssm:eq:sumtail}
 \PP(|Z_M|>u)\le2\exp\left\{-\frac{2M}{M-1}u^2\right\},\quad u\ge0.
\end{equation}
For each fixed $\alpha<2$, $\sup_M\E e^{\alpha Z_M^2}<\infty$. If $a_M$ stays in a bounded set, there exists $\eta>0$ such that
\begin{equation}\label{ssm:eq:shiftUI}
 \sup_M\E e^{(1+\eta)(Z_M+a_M)^2}<\infty.
\end{equation}
\end{corollary}
\begin{proof}
Multiplying \eqref{ssm:eq:mgf} and optimizing the exponential Markov bound gives \eqref{ssm:eq:sumtail}. For $0<\alpha<2$, integration of that tail gives, for example,
\[
 \E e^{\alpha Z_M^2}
 \le1+4\alpha\int_0^\infty u e^{-(2-\alpha)u^2}\,\dd u
 =1+\frac{2\alpha}{2-\alpha}.
\]
Nonpositive $\alpha$ are immediate. Finally, for any $\rho>0$,
\[
 (z+a)^2\le(1+\rho)z^2+(1+\rho^{-1})a^2.
\]
Choose small fixed $\rho,\eta>0$ with $(1+\rho)(1+\eta)<2$ and use boundedness of $a_M$. This proves \eqref{ssm:eq:shiftUI}.
\end{proof}

The gap between the required quadratic coefficient $1$ and the admissible coefficient $2$ is the strict margin that permits the critical number of free coordinates. It also controls third and fourth centered moments at orders $M^{3/2}$ and $M^2$, respectively, by integration of the single-coordinate sub-Gaussian tail.
